\section{The common layer construction}\label{sec:construction}
The three families use the same adjacency rules. We recall them here so that
the seed data in the next sections determine each graph without consulting
another document. A centre will recover the preceding graph as its
neighbourhood; auxiliary vertices will provide the obstructions that make
the next level strict.

\subsection{Seed data and adjacency rules}
Choose a non-empty seed $A$, an integer $d\geq1$, and a graph $R$ on the
type set $\mathcal T=\{0,\ldots,d-1\}$. For each type $j$, choose a
\defi{profile} $P_j\subseteq V(A)$: these are its neighbours in the seed.
Let $\mathcal L_m=\{0,\ldots,m-1\}$, with $\mathcal L_0=\varnothing$.
The graph $G_m$ has vertex set
\[
 V(A)\sqcup\{c_i:i\in\mathcal L_m\}
 \sqcup\{x_{j,i}:j\in\mathcal T,\ i\in\mathcal L_m\}.
\]
Here $c_i$ is the centre of layer $i$ and $x_{j,i}$ is its auxiliary of type $j$.
The following rules specify all edges.
\begin{enumerate}
\item The seed induces $A$.
\item The centres form a clique and are adjacent to every seed vertex.
\item All copies of a fixed auxiliary type form a clique.
\item Auxiliaries of distinct types $j,t$ are adjacent exactly when $jt\in E(R)$.
\item A seed vertex $a$ is adjacent to $x_{j,i}$ exactly when $a\in P_j$.
\item The centre $c_i$ is adjacent to $x_{j,h}$ exactly when $i\ne h$.
\end{enumerate}
Put $B=G_1$, $p=c_0$ and $x_j=x_{j,0}$. Then $G_0=A$, $L_B(p)=A$, and
\begin{equation}\label{eq:general-order}
 |V(G_m)|=|V(A)|+(d+1)m.
\end{equation}
A vertex other than a centre is called \defi{non-central}.

The following two facts supply the positive reductions and will also be used
in the homological arguments.
\begin{lemma}\label{lem:centre}
For $m\geq1$ and $i\in\mathcal L_m$, $L_{G_m}(c_i)\cong G_{m-1}$.
\end{lemma}
\begin{proof}
Its vertices are the seed and all layers except layer $i$. Fix the seed and
renumber the remaining layers in increasing order. All six rules are
preserved, including the inequality of layer indices in the last rule.
\end{proof}

\begin{lemma}\label{lem:remainder}
For $m\geq2$, the graph $G_m-c_i$ belongs to $\D_0$.
\end{lemma}
\begin{proof}
For each type $j$ and $h\ne i$, fold $x_{j,h}$ onto $x_{j,i}$. These
vertices have the same neighbours in the seed and among auxiliary types;
after $c_i$ is deleted, the target sees all remaining centres and the source
omits $c_h$. Hence the required closed-neighbourhood inclusion holds.
The inclusions persist after preceding folds, and every target is retained.
After all folds, each remaining centre is universal. At least one exists
because $m\geq2$, so the resulting cone folds to a vertex.
\end{proof}
For $m=1$ there is no remaining centre. The reduction of $B-p$ must be
checked separately for each seed.

\subsection{Clique counts and the Euler criterion}
Let $\kappa_t(H)$ count cliques of cardinality $t$, with $\kappa_0(H)=1$.
The \defi{clique polynomial} and its alternating sum are
\[
 C_H(z)=\sum_{t\geq0}\kappa_t(H)z^t,\qquad q(H)=C_H(-1).
\]
The \defi{clique complex} $\Cl(H)$ consists of the cliques, including the empty
face. With reduced Euler characteristic, including the empty graph convention,
\begin{equation}\label{eq:euler}
 q(H)=1-\chi(\Cl(H))=-\widetilde\chi(\Cl(H)),\qquad q(\varnothing)=1.
\end{equation}
\begin{lemma}\label{lem:qrec}
For every vertex $v$,
\[
 C_H(z)=C_{H-v}(z)+zC_{L_H(v)}(z),\qquad
 q(H)=q(H-v)-q(L_H(v)).
\]
A cone has $q=0$, and a fold preserves $q$.
\end{lemma}
\begin{proof}
Separate cliques according to whether they contain $v$~\cite{HL1994}.
In a cone, adjoining the apex pairs cliques of opposite signs. For a fold,
Lemma~\ref{lem:dom} makes the neighbourhood term vanish.
\end{proof}
Taking coefficients gives
\begin{equation}\label{eq:clique-coefficients}
 \kappa_t(H)=\kappa_t(H-v)+\kappa_{t-1}(L_H(v))\quad(t\geq1).
\end{equation}
The initial values are $\kappa_0(H)=1$ and $\kappa_t(\varnothing)=0$ for $t\geq1$.

\begin{lemma}\label{lem:qobstruction}
Membership in any $\D_j$, $j\geq-1$, implies $q=0$.
\end{lemma}
\begin{proof}
Start with cones. Inductively, all neighbourhood terms in a level-$j$
reduction vanish. The recurrence preserves $q$ down to $K_1$, where it is zero.
\end{proof}
Thus a non-zero $q$ excludes membership; a zero value does not establish it.

\begin{theorem}\label{thm:template}
Suppose that, for some $s\geq0$,
\begin{enumerate}
\item $A\in\D_s\setminus\D_{s-1}$;
\item $B-p\in\D_s$;
\item $q(L_B(v))\ne0$ for every $v\ne p$.
\end{enumerate}
Then $G_m\in\D_{s+m}\setminus\D_{s+m-1}$ for every $m\geq0$, with
order given by~\eqref{eq:general-order}.
\end{theorem}
This is the Euler criterion used in the article. A proof within this supplement
follows from Theorem~\ref{hom:thm:criterion}: by~\eqref{hom:eq:euler}, a non-zero
$q$ gives non-zero reduced homology, and $\D_s\subseteq\D_{s+1}$ supplies
its remainder hypothesis. Section~\ref{sec:homology} proves that theorem
from the adjacency rules, Lemmas~\ref{lem:centre} and~\ref{lem:remainder},
the first-deletion equivalence~\eqref{eq:first}, and the elementary homology
of Section~\ref{sec:homological-tools}.
No other result of Sections~\ref{sec:family4}--\ref{sec:kernels} is used.

\begin{remark}\label{rem:rest-hypothesis}
The same argument allows $B-p\in\D_{s+1}$ in place of condition (2).
It does not remove the requirement to reduce $B-p$.
\end{remark}

The article also proves that these Euler hypotheses require at least two
auxiliary types, hence at least three new vertices per layer. This restriction
concerns the stated template. It is attained by $F_k$ and is not a lower bound
on the order of every graph of strict level $k$.
